\section{PROOFS}\label{appendix:proofs}

\subsection{Proofs for~\Cref{sec:incoherence}}

\begin{proposition}[Characterisation of $V$ and $Q$]
    For any prior $\pi(a|s)$ and any non-positive reward function $r(a,s)$, we have simple expressions for the soft $Q$ and $V$ functions given by:
    \begin{align*}
        Q^\pi(a_t, s_t) &= \log p_\pi(\OO_{t:T} = 1|s_t,a_t) \\
        V^\pi(s_t) &= \log p_\pi(\OO_{t:T} = 1|s_t)
    \end{align*}
    where we define the auxiliary \emph{optimality variables} $\OO_t$ to have Bernoulli distributions with:
    \[
    p_\pi(\OO_t = 1|s_t,a_t) = e^{r(s_t,a_t)}
    \]
\end{proposition}
\begin{proof}
    We prove this by backward induction. Base case $t = T$, we have:
    \begin{align*}
    Q^\pi(s_T, a_T) &= r(s_T, a_T) = \log (\exp r(s_T, a_T)) \\
    &= \log p_\pi(\OO_{T} = 1|s_T,a_T)
    \end{align*}
    and:
    \begin{align*}
    V^\pi(s_T) &= \log \Expect{a_t \sim \pi(a_T|s_T)}{\exp Q^\pi(s_t, a_T)} \\
    &= \log \Expect{a_t \sim \pi(a_T|s_T)}{p_\pi(\OO_{T} =1|s_T, a_T})\\
    &= \log p_\pi(\OO_{T}=1|s_T)
    \end{align*}
    
    Now, assuming that the hypothesis holds for $t < t' \leq T$, we compute for $t$:
    \begin{align*}
        Q^\pi(&s_{t},a_{t}) = r(s_{t}, a_{t}) + \log \Expect{s_{t+1}}{\exp(V^\pi(s_{t+1})}) \\
        &= \log (p_\pi(\OO_{t} = 1|s_{t},a_{t}) \cdot \Expect{s_{t+1}}{p_\pi(\OO_{t+1:T}|s_{t+1})}) \\
        &= \log \left(p_\pi(\OO_{t} = 1|s_{t},a_{t})  \cdot p_\pi(\OO_{t+1:T}|s_{t+1}, a_{t+1})\right) \\
        &= \log p_\pi(\OO_{t:T} = 1|s_{t},a_{t})
    \end{align*}
    and the proof of the inductive step for $V^\pi(s_t, a_t)$ is the same as the base case shown above.
\end{proof}

\begin{proposition}
    In case of deterministic dynamics, we have that $p(\xi|\OO_{1:T}) = \hat{p}(\xi)$.
\end{proposition}
\begin{proof}
We compute the policy for time steps $t$ and $t+1$:
\begin{align*}
p(a_t|s_t, \OO_{t:T}) &= \frac{p(\OO_{t:T}|a_t, s_t)p(a_t|s_t)}{p(\OO_{t:T}|s_t)} \\
&= \frac{p(\OO_{t:T}|s_{t+1})p(a_t|s_t)}{p(\OO_{t:T}|s_t)}    
\end{align*}
and
\[
p(a_{t+1}|s_{t+1}, \OO_{t+1:T}) = \frac{p(\OO_{t+1:T}|a_{t+1}, s_{t+1})p(a_{t+1}|s_{t+1})}{p(\OO_{t+1:T}|s_{t+1})}
\]
where $\tau(s_{t+1}|s_{t}, a_{t}) = 1$. Thus, considering $\log \hat{p}(\xi)$ we have telescoping series:
\begin{align*}
\log \hat{p}(\xi) &= \sum_{t = 1}^{T} \log p(\OO_{t:T}|s_{t}, a_{t}) - \log p(\OO_{t:T}|s_{t}) + \log p(a_t|s_t) \\
=& \log(\OO_{T:T}|a_T, s_T) - \log p(\OO_{1:T}|s_1) + \sum_{t = 1}^T \log p(a_t|s_t)
\end{align*}
which is exactly:
\[
\log p(\xi|\OO_{1:T}) = \log p(\OO_{1:T}|\xi) - \log(\OO_{1:T}) + \log p(\xi)
\]
\end{proof}
\begin{proposition}
    Given deterministic dynamics $\tau$ and policy $\pi(a|s)$ and an order-respecting $f: \mathbb{R}^{|A|} \to \Delta(|A|)$, such that for $x$ having a unique maximum, $f(x)$ is an $\arg\max$ indicator, $\pi$ is $f$-coherent if and only if it is greedy w.r.t. its own soft‑$Q$ function.
\end{proposition}
\begin{proof}

We first show argmax necessity under order‑respecting $f$: let $x\in\mathbb{R}^{|A|}$ have a unique maximizer $m=\arg\max_i x_i$. If $f$ is order‑respecting and outputs a point mass at $x$, then $f(x)=\delta_m$. Indeed, suppose $f(x)=\delta_i$ with $i\neq m$. Then $x_m>x_i$, but the last condition in the order‑preservation definitoin demands $f_m(x)\ge f_i(x)=1$, a contradiction.

Now, the forward direction. Assume $\pi$ is $f$-coherent, and fix the state $s$. If $Q^\pi(s,\cdot)$ has a unique maximizer, from what we proved above we know that $f(Q^\pi(s,\cdot)) = \delta_{\arg\max}$, i.e. $\pi(s)$ must be an argmax. If there are ties, coherence still forces $\pi(s)$ to be one of the maximizers, which is consistent with the selector’s tie‑break, as otherwise the per‑state $\KL$ contributes $\infty$. So $\pi(s)\in\arg\max Q^\pi(s,\cdot)$.

For the reverse direction, if $\pi(s)\in\arg\max_a Q^\pi(s,a))$ for all $s$, we pick $f$ to be the $\arg\max$ selector $f(Q_\pi(s,\cdot))=\delta_{\pi(s)}$ for all $s$. Hence $\KL(\pi||\pi^{Q,f})=0$ by~\Cref{prop:alternative-incoherence}, i.e., $\pi$ is $f$-coherent.

\end{proof}

% \begin{proposition}[Coherent deterministic policies]
%     If a policy $\pi(a|s)$ and the transition dynamics are deterministic, then $\pi$ is $f$-coherent if and only if $f$ is the $\arg\max$ indicator function.
% \end{proposition}
% \begin{proof}
%     For deterministic dynamics, the equation for soft $Q$ is given by: 
%     \[
%         Q^\pi(s,a) = r(s,a) + V^\pi(s') = r(s, a) + Q^\pi(\tau(s),\pi(\tau(s)))
%     \]
%     where $\tau(s)$ is the unique state such that $\tau(s, a) = \delta_{s'}$ and where $\pi(s)$ is the unique action such that $\pi(a|s) = \delta_{a}$. By induction, this means that:
%     \[
%     Q^\pi(s_t,a_t) = \sum_{t' = t}^T r(s_{t'}, a_{t'})
%     \]
%     for the unique trajectory: \[
%     \xi_{t:T} = (s_t, a_t, s_{t+1}, a_{t+1}, \ldots s_{T}, a_{T}) \sim (\pi, \tau)
%     \]
%     Thus, maximising $Q^\pi(s_0, a_0)$ is equivalent to maximising $\sum_{t} r(s_t, a_t) = J(\pi)$.
% \end{proof}

\begin{proposition}
    The policy $\pi^\BB_T$ is $f$-coherent.
\end{proposition}
\begin{proof}
    We prove by backwards induction that $\pi(s_t|a_t)$ is fixed in iteration $t$ and does not change afterwards. In the base case $t = T$, the $Q$ function $Q^\pi(s_T, \cdot)$ does not depend on $\pi$, so it is fixed in the first iteration and remains unchanged. The inductive case follows because $\pi(s_t|a_t)$ only depends on future times $t' > t$.
\end{proof}

\begin{proposition}
    Given any prior $\pi(a|s)$, there exists a policy $\pi^*(a|s)$ such that we have convergence in distribution:
    \[
        \lim_{\delta \to 0} \pi^\delta \to \pi^*
    \]
    where $\pi^\delta$ is defined with respect to the soft Q function induced by the prior $\pi$. Moreover, the policy $\pi^*$ can be explicitly characterised as the uniform distribution over $A^* := \{a^* \in \A: Q^\pi(s, a^*) = \max_{a \in \A} Q^\pi(s, a)\}$.
\end{proposition}
\begin{proof}
    Let us denote $A_t^* := \{a^* \in \A: Q^\pi(s_t, a^*) = \max_{a \in \A} Q^\pi(s_t, a)\}$ and $a_t^* \in A^*$ arbitrarily chosen. \\
    For any $1 \leq t \leq T$ we have:
    \begin{align*}
    \lim_{\delta \to 0} \pi^\delta(a_t^*|s_t) &= \lim_{\delta \to 0} \frac{\exp\left(\frac{1}{\delta} Q^\pi(s_t, a_t^*)\right)}{\sum_{a \in A} \exp \left(\frac{1}{\delta} Q^\pi(s_t, a)\right)} \\
    &= \lim_{\delta \to 0} \frac{\exp\left(\frac{1}{\delta} Q^\pi(s_t, a^*) - Q^\pi(s_t, a^*)\right)}{\sum_{a \in A} \exp \left(\frac{1}{\delta} Q^\pi(s_t, a) - Q^\pi(s_t, a^*)\right)} \\
    &= \lim_{\delta \to 0} \frac{1}{\sum_{a \in A} \exp \left(\frac{1}{\delta} Q^\pi(s_t, a) - Q^\pi(s_t, a^*)\right)}
    \end{align*}
    where the denominator now splits into:
    \begin{align*}
    &\lim_{\delta \to 0} \sum_{a \in A_t^*} \exp \left(\frac{1}{\delta} Q^\pi(s_t, a) - Q^\pi(s_t, a^*)\right) = |A^*_t|  \\
    &\lim_{\delta \to 0} \sum_{a \in A\backslash A_t^*} \exp \left(\frac{1}{\delta} Q^\pi(s_t, a) - Q^\pi(s_t, a^*)\right) = 0
    \end{align*}
    giving $\pi^\delta(a|s_t)$ uniform on $A_t^*$.
\end{proof}

\begin{corollary}
    A necessary condition for policy $\pi(a|s)$ to be optimal is that:
    \[
    \lim_{\delta \to 0} \kappa_\delta(\pi) < \infty
    \]
\end{corollary}
\begin{proof}
    From~\Cref{def:boltzmann-incoherence},~\Cref{prop:alternative-incoherence} and ~\Cref{lemma:optimal-maxent}, we write:
    \begin{align*}
    \lim_{\delta \to 0} \kappa_\delta(\pi) &= \lim_{\delta \to 0} \KL_\tau(\pi(\tau)|\pi^\delta(\tau)) \\
    &= \lim_{\delta \to 0} \sum_{t = 1}^T \mathbb{E}_{s_t \sim d^t_\pi} \KL(\pi(\cdot|s_t) || \pi^\delta(\cdot|s_t)) \\
    &= \sum_{t = 1}^T \mathbb{E}_{s_t \sim d^t_\pi} \KL(\pi(\cdot|s_t) || \lim_{\delta \to 0} \pi^\delta(\cdot|s_t)) \\
    &= \sum_{t = 1}^T \mathbb{E}_{s_t \sim d^t_\pi} \KL(\pi(\cdot|s_t) || \pi^*(\cdot|s_t))
    \end{align*}
    Thus, for any $s_t \sim d^t_\pi$, the $\KL(\pi|| \pi^*)$ is finite if and only $\text{supp}(\pi^*) \subseteq \text{supp}(\pi)$. Meaning that $\pi$ must assign positive mass to every action in $A_t^*$ for each $t$. But this is a necessary condition for optimality (by backwards induction on $t$).
\end{proof}

\subsection{Proofs for~\Cref{sec:removing-incoherence}}

\begin{proposition}[Strong return improvement lemma]
    The sequence of policies $(\pi^\GG_i)_{i = 0, 1, \dots}$ given by the control-by-inference improves its return monotonically, that is:
    \[
    J(\pi^\GG_{i + 1}) \geq J(\pi^\GG_{i}) 
    \]
\end{proposition}
\begin{proof}    
    We prove this by backwards induction on $T$. First, fix some index $i$ and denote for brevity:
    \[
    \pi = \pi^\GG_i \qquad \pi' = \pi^\GG_{i+1}
    \]
    We also define the future probability of success in a state $s_t$ as in~\Cref{def:softvq}:
    \begin{align*}
    V^\pi(s_t) &= \log p_\pi(\mathcal{O}_{t:T} = 1| s_t) \\
    Q^\pi(s_t, a_t) &= \log  p_\pi(\mathcal{O}_{t:T} = 1| s_t, a_t) = r(s_t, a_t)
    \end{align*}
    Now, we proceed by induction to show that:
    \[
    V^{\pi'}(s_t) \geq V^{\pi}(s_t)
    \]
    The policy $\pi' = \mathcal{G}(\pi)$ can be written from Bayes theorem as:
    \begin{equation}\label{eq:posterior}
    \pi'(a_t|s_t) = \frac{\pi(a_t|s_t) \exp Q^\pi(a_t|s_t)}{\mathbb{E}_{a_t \sim \pi(\cdot|s_t)}[\exp Q^\pi(a|s_t)]}    
    \end{equation}
    Now:
    \begin{equation}\label{eq:inequality}
    \begin{aligned}
        \mathbb{E}_{a_t \sim \pi'(\cdot|s_t)}[\exp Q^\pi(a|s_t)]
        &= \frac{\mathbb{E}_{a_t \sim \pi(\cdot|s_t)}\left[\left(\exp Q^\pi(a|s_t)\right)^2\right]}{\mathbb{E}_{a_t \sim \pi(\cdot|s_t)}[\exp Q^\pi(a|s_t)]} \\
        & \geq \mathbb{E}_{a_t \sim \pi(\cdot|s_t)}[\exp Q^\pi(a|s_t)]
    \end{aligned}
    \end{equation}
    where the first line follows from~\Cref{eq:posterior}, and the second line follows from Cauchy–Schwarz inequality.
    
    First, assume $t=T$. Then, $Q^\pi(a_T|s_T)$ does not depend on $\pi$ and is simply given by $r(s_t, a_t)$. We then immediately have from~\Cref{eq:inequality} that:
    \[
    V^{\pi'}(\mathcal{O}_{T:T}|s_t) \geq V^{\pi}(\mathcal{O}_{T:T}|s_t)
    \]
    establishing the base of the induction. For the inductive case $t < T$, assumethat the above holds for all $t < t' \leq T$ and write:
    \begin{align*}
      \exp V^{\pi'}(s_t) &= p_{\pi'}(\mathcal{O}_{t:T}=1 | s_t)\\
      &= \mathbb{E}_{a_t\sim \pi'(\cdot| s_t)}\left[e^{r(s_t,a_t)}
      \mathbb{E}_{s_{t+1}} \exp V^{\pi'}(s_{t+1})\right] \\
      &\geq \mathbb{E}_{a_t\sim \pi'(\cdot| s_t)}\left[e^{r(s_t,a_t)}
      \mathbb{E}_{s_{t+1}} \exp V^{\pi}(s_{t+1})\right] \\
      &= \mathbb{E}_{a_t\sim \pi'(\cdot| s_t)}\left[\exp Q^\pi(s_t, a_t)\right] \\
      &\geq \mathbb{E}_{a_t\sim \pi(\cdot| s_t)}\left[\exp Q^\pi(s_t, a_t)\right] \\
      &= \exp V^\pi(s_t)
    \end{align*}
    where the first inequality comes from the inductive hypothesis, and the second from~\Cref{eq:inequality}.
    From the Cauchy–Schwarz, inequalities are strict unless $Q^\pi(s_t,\cdot)$ is $\pi(\cdot| s_t)$-almost-surely constant. Intuitively, the update is non‑trivial exactly when the old policy is not already proportional to the frozen‑future success factors.
\end{proof}

\subsection{Proof of~\Cref{sec:equivalence}}

We split the proof of~\Cref{theorem:equivalence-thrice} into several lemmas.

\begin{lemma}~\label{lemma:equivalence-temperature-folding}
    For any prior $p(a|s)$, non-positive reward function $r(s_t,a_t)$, deterministic dynamics $\tau(s_{t+1}|s_t, a_t)$, and $k \in \mathbb{N}_+$ we have $\pi_{\alpha(k)} = \pi^\HH_k$, for $\alpha(k) = 2^k$.
\end{lemma}
\begin{proof}
    Throughout the proof, we will assume that the prior is uniform; otherwise the prior is folded into the reward. Let us first consider the case of $T = 1$. We have that:
    \begin{align*}
        \log p(a|s, \mathcal{O}^{(k)}) &= \log\left(\frac{p(\mathcal{O}^{(k)}|a, s)p(a|s)}{p(\mathcal{O}^{(k)}|s)}\right) \\
        &= r_k(s, a) + C
    \end{align*}
    where the constant
    \[
    C = \log |A| - \log \frac{\sum_{a}\exp r_k(s, a)}{|A|}
    \]
    does not depend on $a$. We also observe that shifting reward by any additive constant does not change the policy (a soft version of the reward shaping lemma): for $r'(s, a) = r(s,a) + C$ we have
    \begin{align*}
        p(a|s, \mathcal{O}') &= \frac{p(a|s) e^{r(s, a) + C}}{\sum_{a'} p(a'|s)e^{r(s,a) + C}} \\
        &= \frac{p(a|s) e^{r(s, a)}}{\sum_{a'} p(a'|s)e^{r(s,a)}}
    \end{align*}
    To ignore the constants, we write $r \equiv r'$ when $r = r' + C$. Thus, the iterative process of adding the posterior to the reward results in a sequence:
    \begin{align*}
    r_0(s, a) &\equiv r(s, a) \\
    r_1(s, a) &\equiv r_0(s, a) + \log p(a|s, \mathcal{O}^{0}) \\
    &\equiv 2r(s, a) + C_1 \equiv 2r(s,a) \\
    % r_2(s, a) &\equiv r_1(s, a) + \log p(a|s, \mathcal{O}^{(2)}) \\
    % &\equiv 2r_1(s, a) + C_2 \equiv 4r(s,a)\\
    \dots\\
    r_{k}(s, a) &\equiv r_{k-1}(s, a) + \log p(a|s, \mathcal{O}^{(k-1)}) \\
    &\equiv 2r_{k-1}(s, a) + C_k \equiv 2^{k}r(s,a)
    \end{align*}
    This proves, by induction on $k$, that $\pi^\HH_k = \pi_{\alpha(k)}$.\\
    The case of $T > 1$ is then done by backwards induction over the time horizon. We are given $1 \leq t < T$ and we assume that the following inductive hypothesis holds for all $t < t'$ and for all $k$:
    \[
    r_k(s, a) = 2^kr_0(s, a) \qquad p(\mathcal{O}_{t':T}^{\alpha(k)}|s_{t'}) = p(\mathcal{O}_{t':T}^{(k)}|s_{t'})
    \]
    In case of the policy obtained by decreasing the temperature, we have:
    \begin{align*}
        \pi_{\alpha(k+1)}(a_t|s_t) &= \log p(a_t|s_t, \mathcal{O}_{t:T}^{\alpha(k+1)}) \\
        & \equiv \log p(\mathcal{O}_{t:T}^{\alpha(k+1)}|a_t, s_t) + \log p(a_t|s_t) - \log p(\OO_{t:T}|s_t)\\
        & \equiv \log p(\mathcal{O}_{t}^{\alpha(k+1)}|a_t, s_t) + \log p(\mathcal{O}_{t+1:T}^{\alpha(k+1)}|a_t, s_t) \\
        &= 2^{k + 1}r_0(a_t, s_t) + \log p(\mathcal{O}_{t+1:T}^{\alpha(k+1)}|s_{t+1})
    \end{align*}
    where the first line follows from the definition, the next one from the Bayes' law, the next one from the assumption that the prior is uniform, and the next one from the definition of the $\mathcal{O}_{1}^{\alpha(k+1)}$ and using the fact that there is a unique $s_{t+1}$ which follows $\tau(a_t, s_t)$.
    
    We look at the policy obtained by folding the posterior into the reward:
    \begin{align*}
        \pi_{(k+1)}^\HH(a_t|s_t) &= \log p(a_t|s_t, \mathcal{O}_{t:T}^{(k+1)}) \\
        & \equiv \log p(\mathcal{O}_{t}^{(k+1)}|a_t, s_t) + \log p(\mathcal{O}_{t+1:T}^{(k+1)}|a_t, s_t) \\
        &= r_{k+1}(a_t, s_t) + \log p(\mathcal{O}_{t+1:T}^{(k+1)}|s_{t+1}) \\
        &= 2^{k+1}r_0(a_t, s_t) + \log p(\mathcal{O}_{t+1:T}^{\alpha(k+1)}|s_{t+1})
    \end{align*}
    where the last line follows from the inductive hypothesis.
\end{proof}

\begin{lemma}~\label{lemma:equivalence-folding-retraining}
    For any prior $p(a|s)$, non-positive reward function $r(s_t,a_t)$, arbitrary dynamics $\tau(s_{t+1}|s_t, a_t)$, and $k \in \mathbb{N}_+$ we have $\pi^\mathcal{G}_{k} = \pi^\FF_k$.
\end{lemma}
\begin{proof}
Proof by induction on $k$. Assume without loss of generality uniform prior $p$, otherwise fold it into the reward. For $k=0$ they coincide by definition. Now, assume that $\pi_k^\GG = \pi_k^\FF$, and we show the inductive step for $k+1$. Let us denote the corresponding $Q$ functions by $Q^\GG_k$ and $Q^\FF_k$.
First, exactly as in the proof of~\Cref{prop:strong-improvement}, we know that:
\[
        \pi_{k+1}^\GG(a_t|s_t) = \frac{\pi_{k}^\GG(a_t|s_t) \exp Q^{\GG}_k(a_t|s_t)}{\mathbb{E}_{a_t \sim \pi(\cdot|s_t)}[\exp Q^{\GG}_k(a|s_t)]}\ \propto\ \pi_{k}^\GG(a_t|s_t) \exp Q^{\GG}(a_t|s_t)
\]
and on the other hand, we also know from the definition that:
\[
    \pi_{k+1}^\FF(a_t|s_t) \ \propto \ p(s_t|a_t) \exp Q^{\FF}_k(a_t|s_t)
\]
Thus, we aim to show that: \[
p(s_t|a_t) \exp Q^{\FF}_k(a_t|s_t) = C_t\pi_{k}^\GG(a_t|s_t) \exp Q^{\GG}_k(a_t|s_t)
\]
up to some multiplicative constant $C_t$. We do this again by induction, this time on a time horizon $t$.
Using~\Cref{def:softvq}, we have:
\[
  \exp Q^{\FF}_k(s_T,a_T)=\exp r_k(s_T,a_T)
  =\pi^\GG_k(a_T| s_T)\exp r_0(s_T,a_T)
  =\pi^\GG_k(a_T| s_T)\exp Q^{\GG}_k(s_T,a_T)
\]
which proves this in case $t=T$. Inductive step: assume that the hypothesis holds for $t+1$, and write:
\[
\exp V^\FF(s_{t+1})
=\mathbb{E}_{a_{t+1}\sim p}[\exp Q^{\FF}_k(s_{t+1}(a_{t+1})]
=\frac{C_t}{|A|}\mathbb{E}_{a_{t+1}\sim \pi^\GG_k}[\exp Q^\GG_k(s_{t+1}(a_{t+1})]
\]
Taking $\log$ and bringing the factor outside the expectation shows that:
 \[
  \mathbb{E}_{s_{t+1}}[\exp V^\FF_k(s_{t+1})]
  =\frac{C_t}{|A|}\mathbb{E}_{s_{t+1}}[\exp V^\GG_k(s_{t+1})]
\]
which means:
\[
\exp Q^\FF_k(s_t,a_t)
  =\exp r_k(s_t,a_t)\cdot \frac{C_t}{|A|}
  \mathbb{E}_{s_{t+1}}[\exp V^\GG_k(s_{t+1})]
\]
proving the lemma.
\end{proof}

\begin{lemma}~\label{lemma:equivalence-folding-orig}
    For any prior $p(a|s)$, non-positive reward function $r(s_t,a_t)$, deterministic dynamics $\tau(s_{t+1}|s_t, a_t)$, and $k \in \mathbb{N}_+$ we have $\pi^\FF_{k} = \pi_{\alpha(k)}$ for schedule $\alpha(k) = k$.
\end{lemma}
\begin{proof}
 The proof follows almost exactly like the one in~\Cref{lemma:equivalence-temperature-folding}. The base case is identical. In the inductive case, we write analogously:
\begin{align*}
    r_0(s, a) &\equiv r(s, a) \\
    r_1(s, a) &\equiv r_0(s, a) + \log p(a|s, \mathcal{O}^{0}) \\
    &\equiv 2r(s, a) + C_1 \equiv 2r(s,a) \\
    % r_2(s, a) &\equiv r_1(s, a) + \log p(a|s, \mathcal{O}^{(2)}) \\
    % &\equiv 2r_1(s, a) + C_2 \equiv 4r(s,a)\\
    \dots\\
    r_{k}(s, a) &\equiv r_0(s, a) + \log p(a|s, \mathcal{O}^{(k-1)}) \\
    &\equiv kr(s, a) + C_k \equiv kr(s,a)
    \end{align*}
\end{proof}
and the rest of the proof follows as before.


\subsection{Proof of~\Cref{corr:rate-of-convergence}}

\begin{lemma}~\label{lemma:argmax-derivative}
Given a function $h(x) = f(\arg\max_t (f(t) + x g(t)))$ for some differentiable functions $f, g$, the derivative of $h(x)$ is:
    \[
        \frac{dh}{dx} = -\frac{f'(t^*(x)) g'(t^*(x))}{f''(t^*(x)) + x g''(t^*(x))}
    \]
    where $t^*(x) = \arg\max_t f(t) + x g(t)$.
\end{lemma}
\begin{proof}
Let $t^*(x) = \arg\max_t (f(t) + x g(t))$. Therefore, we have:
\[
    \frac{d}{dt} \left( f(t) + x g(t) \right) \bigg|_{t = t^*(x)} = 0 = f'(t^*(x)) + x g'(t^*(x))
\]
Thus:
\[ x = -\frac{f'(t^*(x))}{g'(t^*(x))} \]

To differentiate \( h(x) \), we use the fact that \( h(x) = f(t^*(x)) \). This gives us:
\[ \frac{dh}{dx} = \frac{d}{dx} f(t^*(x)) \]

By the chain rule:
\[ \frac{dh}{dx} = f'(t^*(x)) \cdot \frac{dt^*(x)}{dx} \]

To find \( \frac{dt^*(x)}{dx} \), we differentiate the first-order condition \( f'(t^*(x)) + x g'(t^*(x)) = 0 \) with respect to \( x \):
\[ \frac{d}{dx} \left( f'(t^*(x)) + x g'(t^*(x)) \right) = 0 \]

Applying the chain rule:
\[ f''(t^*(x)) \cdot \frac{dt^*(x)}{dx} + g'(t^*(x)) + x g''(t^*(x)) \cdot \frac{dt^*(x)}{dx} = 0 \]

Rearrange to solve for \( \frac{dt^*(x)}{dx} \):
\[ \left( f''(t^*(x)) + x g''(t^*(x)) \right) \frac{dt^*(x)}{dx} = -g'(t^*(x)) \]
\[ \frac{dt^*(x)}{dx} = -\frac{g'(t^*(x))}{f''(t^*(x)) + x g''(t^*(x))} \]

Finally, substituting this back into the expression for \( \frac{dh}{dx} \):
\[ \frac{dh}{dx} = f'(t^*(x)) \cdot \left( -\frac{g'(t^*(x))}{f''(t^*(x)) + x g''(t^*(x))} \right) \]

\end{proof}

\begin{corollary}[Return improvement rate]
    In case of deterministic dynamics $\tau$, given a sequence of policies $\pi^\GG_i$, for $k \in \mathbb{N}_+$ have that:
    \[
    J(\pi^\GG_{k}) - J(\pi^\GG_{k-1}) = \frac{1}{{k}} \cdot \frac{J'(\pi^\GG_{k}) \cdot \hat{\HH}'(\pi^\GG_{k})}{J''(\pi^\GG_k) + \frac{1}{{k}} \hat{\HH}''(\pi^\GG_{k})} + O\left(\frac{1}{{k}^2}\right)
    \]
    where $\hat{\HH}(\pi)$ denotes the causal entropy of policy $\pi$, and the derivatives are taken with respect to the temperature $\alpha$.
\end{corollary}
% \jacek{The proof and the bounds follow from the equivalence to temperature sequence. Bounds should be something like $\alpha \HH(\pi_i)$.}
\begin{proof}
    To prove this, we use an alternative characterization of the maximum entropy  policies~\citep{haarnoja_reinforcement_2017,levine_reinforcement_2018}. For a fixed temperature $\alpha \in \RR_+$, the policy $\pi_\alpha$ maximizes the functional:
    \[
    J_\alpha(\pi) = \mathbb{E}_{s_t, a_t}[r(s_t, a_t) + \frac{1}{\alpha} \mathcal{H}(\pi(\cdot|s_t))]= J(\pi) + \frac{1}{\alpha} \hat{\HH}(\pi)
    \]
    where the expectation is taken over the trajectory induced by the policy $\pi$, $\HH(\pi(\cdot|s))$ denotes the Shannon entropy of the policy in state $s$, and $\hat{\HH}$ denotes the causal entropy of the policy $\pi$. Using the characterisation of the iterative retraining given by~\Cref{theorem:equivalence-thrice} we know that $\pi^\GG_j = \pi_{\alpha(j)}$ for $\alpha(j) = j$.

    From the Taylor expansion, for any twice-differentiable $u(\alpha)$ we have:
    \[
    u(\alpha + \eta) - u(\alpha) = \eta u'(\alpha) + O(\eta^2)
    \]
    Substituting $u(\alpha) = J(\arg\max_\pi J_{\frac{1}{\alpha}}(\pi))$, we obtain: 
    \begin{align*}
        J(\pi^\GG_{k-1}) - J(\pi^\GG_{k})
        &= J(\pi_{\alpha(k - 1)}) - J(\pi_{\alpha(k)}) \\
        &= u\left(\frac{1}{{k}} + \eta\right) - u\left(\frac{1}{{k}}\right)  \\
        &= \eta u'\left(\frac{1}{{k}}\right) + O\left(\eta^2\right)
    \end{align*}
    for $\eta = \frac{1}{k({k-1})}$. Now, using~\Cref{lemma:argmax-derivative}, we derive:
    \[
    u'(x) = -\frac{J'(\pi^*(x)) \hat{\HH}'(\pi^*(x))}{J''(\pi^*(x)) + x \hat{\HH}''(\pi^*(x))}
    \]
    and substituting $x = \frac{1}{{k}}$, and therefore $\pi^*(x) = \pi_{\alpha(k)} = \pi^\GG_k$, we get:
    \[
    u'(x) = -\frac{J'(\pi_{\alpha(k)}) \hat{\HH}'(\pi_{\alpha(k)})}{J''(\pi^\GG_k) + \frac{1}{{k}} \hat{\HH}''(\pi_{\alpha(k)})}
    \]
    which gives us the final equation:
    \[
    J(\pi^\GG_{k}) - J(\pi^\GG_{k-1}) = \frac{1}{{k}} \frac{J'(\pi^\GG_{k}) \hat{\HH}'(\pi^\GG_{k})}{J''(\pi^\GG_{k})) + \frac{1}{{k}} \hat{\HH}''(\pi^\GG_{k})} + O\left(\frac{1}{k}^2\right)
    \]
\end{proof}
